\chapterbackmatter{Solutions to Supplementary Exercises}

\begin{enumerate}
  \SupplementarySolution{1}{Axion--Kalb--Ramond Duality from a Parent Action}
  It is cleanest to use the differential-form normalization stated in the
  problem; numerical factors for the component normalization follow by
  rescaling \(H\) and \(a\).
  \begin{enumerate}
    \item Variation with respect to \(a\) gives
      \[
        dH=0,
        \qquad\text{or equivalently}\qquad
        \epsilon^{\mu\nu\rho\sigma}\partial_\mu H_{\nu\rho\sigma}=0.
      \]
      On a contractible patch, the Poincare lemma implies
      \[
        H=dB,
        \qquad H_{\mu\nu\rho}=3\partial_{[\mu}B_{\nu\rho]}.
      \]
      The action becomes
      \[
        S_B=-\frac1{2g^2}\int dB\wedge *dB,
      \]
      and is invariant under \(B\mapsto B+d\Lambda_1\).  Global closed
      but non-exact sectors of \(H\) must be summed separately when the
      spacetime has nontrivial third cohomology.

    \item Up to a boundary term,
      \[
        \int a\,dH=-\int da\wedge H.
      \]
      Varying the independent three-form then gives
      \[
        *H=g^2 da,
        \qquad H=g^2*da
      \]
      with the displayed orientation convention.  Substitution completes
      the square and yields
      \[
        S_a=-\frac{g^2}{2}\int da\wedge *da.
      \]
      Thus a kinetic coefficient proportional to \(g^{-2}\) on the
      two-form side becomes one proportional to \(g^2\) on the scalar
      side: the dual coupling is inverted.  The two descriptions each
      carry one local degree of freedom in four dimensions.

    \item For a \(p\)-form potential \(A_p\), take an independent
      \((p+1)\)-form \(H\) and a dual potential
      \(\widetilde A_{D-p-2}\), with
      \(\widetilde F=d\widetilde A\).  A parent action is
      \[
        S[H,\widetilde A]
          =-\frac1{2g^2}\int H\wedge *H
           +\int H\wedge\widetilde F,
      \]
      where an orientation-dependent sign may be placed in the second
      term.  Eliminating \(\widetilde A\) gives \(dH=0\), hence
      \(H=dA\).  Eliminating \(H\) instead gives
      \[
        \widetilde F=\pm g^{-2}*H
      \]
      and a dual kinetic coefficient proportional to \(g^2\).  The ranks
      are
      \[
        \deg H=p+1,\qquad
        \deg\widetilde F=D-p-1,\qquad
        \deg\widetilde A=D-p-2.
      \]
  \end{enumerate}

  \SupplementarySolution{2}{Massive Antisymmetric-Tensor Duality}
  \begin{enumerate}
    \item With \(F_{\mu_0\cdots\mu_q}=(q+1)
      \partial_{[\mu_0}A_{\mu_1\cdots\mu_q]}\), the mostly-plus
      Lorentzian Lagrangian is
      \[
        \mathcal L_A
          =-\frac1{2(q+1)!}F_{\mu_0\cdots\mu_q}F^{\mu_0\cdots\mu_q}
           -\frac{m^2}{2q!}A_{\mu_1\cdots\mu_q}A^{\mu_1\cdots\mu_q}.
      \]
      Its equation of motion is \(d^\dagger F-m^2A=0\).  Applying
      \(d^\dagger\) gives \(d^\dagger A=0\) for \(m\ne0\), the massive
      transversality condition.

    \item The algebra is most transparent after Euclidean continuation.
      Let \((X,Y)=\int X\wedge *Y\), and let
      \(B\) have rank \(D-q-1\).  A Gaussian parent action is
      \[
        S_{\rm P}^{E}[A,B]
          =\frac12(B,B)-i\int B\wedge dA
            +\frac{m^2}{2}(A,A),
      \]
      with the harmless phase and a rank-dependent sign fixed by the
      convention for \(*^2\).  The \(B\) equation says
      \[
        B=\pm i*dA.
      \]
      Substitution gives the Euclidean continuation of the original
      kinetic term plus \(m^2(A,A)/2\).

      Alternatively, integrate the mixed term by parts.  The \(A\)
      equation is algebraic and gives
      \[
        m^2A=\pm i*dB.
      \]
      Substituting it and defining \(\widetilde A=B/m\) gives, after
      returning to Lorentzian signature,
      \[
        \mathcal L_{\widetilde A}
          =-\frac1{2(D-q)!}(d\widetilde A)^2
           -\frac{m^2}{2(D-q-1)!}\widetilde A^2 .
      \]
      Hence the dual potential has rank \(D-q-1\) and the same mass.  Up
      to the orientation signs, the first-order duality equations are
      \[
        m\widetilde A=*dA,
        \qquad d\widetilde A=m*A.
      \]

    \item A massive particle in \(D\) dimensions has little group
      \(SO(D-1)\).  The original field is a rank-\(q\) antisymmetric tensor
      of that group, so
      \[
        N_A=\binom{D-1}{q}.
      \]
      The dual field has rank \(D-q-1\), and therefore
      \[
        N_{\widetilde A}
          =\binom{D-1}{D-q-1}
          =\binom{D-1}{q}=N_A.
      \]

    \item The identity underlying the construction is the functional
      Gaussian transform
      \[
        e^{-\frac12(*dA,*dA)}
          =\mathcal{N}\int\mathcal{D}B\,
             \exp\!\left[-\frac12(B,B)+i\int B\wedge dA\right],
      \]
      with gauge zero modes treated separately.  Integrating \(B\) is the
      inverse transform; integrating \(A\) produces the dual description.

      At \(m=0\), the step \(A=\pm i*dB/m^2\) and the rescaling
      \(B=m\widetilde A\) are singular.  The parent integral itself still
      makes sense: integration over \(A\) now imposes \(dB=0\).  Locally
      \(B=dC\), where \(C\) has rank \(D-q-2\), the familiar massless
      dual rank.  The count changes consistently to
      \[
        \binom{D-2}{q}
          =\binom{D-2}{D-q-2},
      \]
      because both massless potentials have gauge redundancy.
  \end{enumerate}

  \SupplementarySolution{3}{Graviton and Gravitino Counts in Even Dimensions}
  \begin{enumerate}
    \item A massless particle in \(D\) dimensions transforms under the
      compact little group \(SO(D-2)\).  Put \(n=D-2\).  The graviton is a
      symmetric traceless rank-two tensor of this group.  A symmetric
      tensor has \(n(n+1)/2\) components, and removal of its trace gives
      \[
        N_g=\frac{n(n+1)}2-1
          =\boxed{\frac{D(D-3)}2}.
      \]
    \item Let \(s_{D-2}\) denote the real physical dimension of the
      appropriate irreducible little-group spinor after the required
      reality and chirality conditions have been imposed.  A vector-spinor
      initially has \((D-2)s_{D-2}\) components.  Its gamma trace is an
      arbitrary spinor with \(s_{D-2}\) components, so the gamma-traceless
      gravitino has
      \[
        N_{\psi_\mu}
          =\bigl((D-2)-1\bigr)s_{D-2}
          =\boxed{(D-3)s_{D-2}}
      \]
      physical polarizations.
    \item For \(D=10\), the little group is \(SO(8)\), and an irreducible
      chiral spinor has \(s_8=8\).  Thus
      \[
        N_g=35,\qquad N_{\psi_\mu}=7\cdot8=56.
      \]
      These two fields alone do not have equal bosonic and fermionic
      counts.  The minimal ten-dimensional supergravity multiplet also
      contains a two-form \(B_{\mu\nu}\), a dilaton \(\phi\), and a
      chiral dilatino \(\chi\).  Their little-group counts are
      \[
        N_B=\binom82=28,\qquad N_\phi=1,\qquad N_\chi=8.
      \]
      Consequently
      \[
        N_{\rm bos}=35+28+1=64,
        \qquad
        N_{\rm ferm}=56+8=64.
      \]
      The two-form and dilaton therefore complete the bosonic sector,
      while the dilatino completes the fermionic sector.
  \end{enumerate}

  \SupplementarySolution{4}{A Spin-Three Supersymmetry Bound}
Let \(q\) be the number of real supercharges.  After dimensional
  reduction to four dimensions, a massless representation has \(q/4\)
  fermionic creation operators.  Each lowers helicity by \(1/2\), so the
  full multiplet has helicity span
  \[
    \Delta h=\frac12\frac q4=\frac q8.
  \]
  A CPT-self-conjugate placement minimizes the largest absolute
  helicity by centering this range about zero.  The smallest possible
  value of the maximum spin is therefore
  \[
    j_{\rm min}=\frac{\Delta h}{2}=\frac q{16}.
  \]
  Requiring every state to have spin at most three gives
  \[
    q\leq48.
  \]

  A minimal real spinor in eleven-dimensional Minkowski space has
  \(32\) real components, so eleven-dimensional supersymmetry satisfies
  this bound.  In twelve dimensions the relevant minimal real
  supersymmetry has \(64\) real components.  Its four-dimensional
  reduction would require
  \(j_{\rm min}=64/16=4\), already above the allowed ceiling.  Minimal
  spinor dimensions do not decrease in higher dimensions.  Hence
  \[
    \boxed{D_{\max}=11}.
  \]
  Allowing spin three instead of spin two enlarges the abstract
  supercharge bound from \(32\) to \(48\), but the next available
  higher-dimensional spinor has \(64\) components; there is no
  intermediate spacetime dimension that realizes the extra allowance.

  \SupplementarySolution{5}{Type-II and Eleven-Dimensional Reductions}
  \begin{enumerate}
    \item In nine dimensions the little group is \(SO(7)\).  The useful
      polarization counts are
      \[
        g_{\mu\nu}:27,\quad A_\mu:7,\quad
        B_{\mu\nu}:21,\quad C_{\mu\nu\rho}:35,\quad
        \chi:8,\quad \psi_\mu:48.
      \]
      Type IIA has ten-dimensional bosons
      \(g,B_2,\phi,C_1,C_3\).  On a circle they decompose as
      \[
      \begin{array}{c@{\quad\longrightarrow\quad}l}
        g&g_9+A_1+\varphi,\\
        B_2&B_2+A_1,\\
        \phi&\varphi,\\
        C_1&A_1+\varphi,\\
        C_3&C_3+B_2.
      \end{array}
      \]
      Hence the nine-dimensional bosonic content is
      \[
        g_9+3A_1+3\varphi+2B_2+C_3
      \]
      and has
      \[
        27+3(7)+3+2(21)+35=128
      \]
      polarizations.

      Type IIB has \(g\), two two-forms, two scalars, and a self-dual
      four-form.  The metric contributes \(g_9+A_1+\varphi\), each
      two-form contributes \(B_2+A_1\), and the two scalars remain
      scalars.  The circle reduction of the self-dual four-form supplies
      one independent nine-dimensional three-form.  Thus IIB gives the
      same list
      \(g_9+3A_1+3\varphi+2B_2+C_3\).

      In either theory the fermions become two nine-dimensional
      gravitini and four spinors:
      \[
        2(48)+4(8)=128.
      \]
      Nine-dimensional spinors have no chirality label analogous to the
      ten-dimensional one.  The opposite-chirality IIA and
      same-chirality IIB assignments therefore reduce to the same
      nonchiral maximal nine-dimensional multiplet.
    \item Reducing eleven-dimensional supergravity on a two-torus gives
      the same result directly.  The metric polarizations branch as
      \[
        44=27+2(7)+3,
      \]
      corresponding to a nine-dimensional graviton, two vectors, and
      the three components of the internal two-torus metric.  The
      three-form branches as
      \[
        84=35+2(21)+7,
      \]
      corresponding to a three-form, two two-forms, and one vector.
      Together these are precisely
      \(g_9+3A_1+3\varphi+2B_2+C_3\), with \(128\) bosonic states.
      The eleven-dimensional gravitino branches as
      \[
        128=2(48)+4(8),
      \]
      reproducing the same \(128\) fermionic states.
    \item Compactification of type IIB on a six-torus preserves all
      \(128+128\) massless zero-mode polarizations and gives maximal
      \({\cal N}=8\) supersymmetry in four dimensions.  Its fields can be
      organized as
      \[
      \begin{array}{c|c|c}
        \text{field}&\text{multiplicity}&\text{physical states}\\ \hline
        g_{\mu\nu}&1&2\\
        A_\mu&28&56\\
        \varphi&70&70\\ \hline
        \psi_\mu&8&16\\
        \chi&56&112
      \end{array}
      \]
      so that
      \[
        N_{\rm bos}=2+56+70=128,\qquad
        N_{\rm ferm}=16+112=128.
      \]
      The lower-dimensional multiplet is again nonchiral and
      CPT-complete.
  \end{enumerate}

  \SupplementarySolution{6}{Kaluza--Klein BPS Towers and the Eleventh Dimension}
  \begin{enumerate}
    \item \emph{Kaluza--Klein Momentum as Central Charge.}
    Split the eleven-dimensional index as \(M=(\mu,10)\).  The momentum
      term in the eleven-dimensional algebra is
      \[
        \{Q,Q^{\mathsf T}\}\supset i\Gamma^M C P_M
          =i\Gamma^\mu C P_\mu+i\Gamma^{10}C P_{10}.
      \]
      Ten-dimensional Lorentz transformations act only on \(\mu\).
      Therefore \(P_{10}\) has no ten-dimensional Lorentz index and is a
      scalar.  Translation invariance along the circle makes it commute
      with all ten-dimensional Poincare generators.

      Under the decomposition into two ten-dimensional Majorana-Weyl
      supercharges of opposite chirality, \(\Gamma^{10}\) is the chirality
      operator.  The second term is consequently the scalar central-charge
      entry of the type-IIA algebra.  The geometric origin also shows why
      this central charge is conserved: it is ordinary momentum in the
      higher-dimensional theory.
    \item \emph{The Kaluza--Klein BPS Tower.}
    Periodicity around a circle of radius \(R\) quantizes momentum:
      \[
        P_{10}=\frac nR,\qquad n\in\mathbb Z.
      \]
      A state massless in eleven dimensions obeys
      \(P_M P^M=0\), so from the ten-dimensional viewpoint
      \[
        m_n^2=P_{10}^2,\qquad
        \boxed{m_n=\frac{|n|}{R}}.
      \]
      Since \(Z=P_{10}\), this saturates \(m\geq|Z|\).

      Saturation makes sixteen of the thirty-two real supercharges null.
      The other sixteen form eight creation-annihilation pairs, giving
      \(2^8=256\) states.  At each nonzero \(n\), these transform as
      \({\bf44}\oplus{\bf84}\) bosons and \({\bf128}\) fermions of the
      massive little group \(SO(9)\).  The levels \(n\) and \(-n\) carry
      opposite central charge and are exchanged by charge conjugation.
    \item \emph{The Emerging Eleventh Dimension.}
    The first Kaluza--Klein excitation has mass \(m_{\rm KK}=1/R_{11}\).
      Equating it to the D0-brane mass gives
      \[
        \frac1{R_{11}}=\frac1{g_s\ell_s},
        \qquad\boxed{R_{11}=g_s\ell_s}.
      \]
      The \(n\)-D0 threshold states then have masses
      \(|n|/(g_s\ell_s)=|n|/R_{11}\), exactly the Kaluza--Klein tower.

      As \(g_s\to\infty\) at fixed string length, \(R_{11}\) grows and the
      spacing \(1/R_{11}\) collapses.  The tower that was heavy in weakly
      coupled ten-dimensional string theory becomes the continuum of
      momentum in a decompactifying eleventh dimension.
  \end{enumerate}

  \SupplementarySolution{7}{Type-II Chirality and Dimensional Reduction}
  \begin{enumerate}
    \item \emph{IIA and IIB Chirality.}
    In a representation adapted to ten dimensions, the
      eleven-dimensional gamma matrices consist of the ten
      \(\Gamma^\mu\) plus the ten-dimensional chirality matrix
      \(\Gamma^{10}=\Gamma_*\).  Projecting an eleven-dimensional Majorana
      spinor with
      \[
        Q_\pm=\frac12(1\pm\Gamma_*)Q
      \]
      gives two real sixteen-component ten-dimensional Weyl spinors.
      Because the two projectors have opposite eigenvalues, the result is
      the nonchiral type-IIA pair.

      Type IIB instead has two real Weyl supercharges of the same
      ten-dimensional chirality.  An ordinary circle reduction cannot
      change the algebraic fact that a single eleven-dimensional spinor
      decomposes into opposite chiralities.  Thus it yields IIA, not IIB.
      IIB is related after a further compactification by T-duality, which
      reverses one right-moving spacetime chirality.
    \item \emph{Reducing the Eleven-Dimensional Three-Form.}
    Write the eleven-dimensional potential as
      \[
        A_{MNP}\longrightarrow
          A_{\mu\nu\rho},\qquad
          B_{\mu\nu}=A_{\mu\nu\,10}.
      \]
      These are the type-IIA Ramond--Ramond three-form and
      Neveu--Schwarz two-form.  Components with two circle indices vanish
      by antisymmetry.

      The same index split applies to the algebraic charges:
      \[
        Z_{MN}\longrightarrow Z_{\mu\nu},\ Z_{\mu\,10},
      \]
      \[
        Z_{M_1\cdots M_5}\longrightarrow
          Z_{\mu_1\cdots\mu_5},\
          Z_{\mu_1\cdots\mu_4\,10}.
      \]
      Together with \(P_{10}\), these become the scalar, one-, two-,
      four-, and five-form structures allowed by the IIA algebra, subject
      to the ten-dimensional chirality and Hodge-duality relations.
    \item \emph{Nine-Dimensional Unification.}
    A ten-dimensional Weyl spinor has sixteen real components.  After
      restriction to \(\operatorname{Spin}(1,8)\), the two
      ten-dimensional chiral spinor
      representations both become the same sixteen-component
      nine-dimensional Majorana representation: nine dimensions has no
      chirality operator independent of the gamma matrices.

      Thus both reductions contain two nine-dimensional Majorana
      supercharges.  A sign or interchange in the second supercharge,
      accompanied by an \(R\)-symmetry basis change among the scalar and
      vector charges, maps the reduced IIA algebra to the reduced IIB
      algebra.  T-duality realizes precisely this map: it exchanges circle
      momentum and winding and flips the chirality convention of one
      ten-dimensional supercharge.  The equality of the nine-dimensional
      algebras is therefore the representation-theoretic shadow of
      IIA/IIB T-duality.
  \end{enumerate}

  \SupplementarySolution{8}{Tensorial Charges and the \(528\)-Component M-Algebra}
  \begin{enumerate}
    \item \emph{Topological \(p\)-Brane Charge.}
    On a spatial worldvolume \(\Sigma_p\), define
      \[
        Z^{\mu_1\cdots\mu_p}
          =T_p\int_{\Sigma_p}
            dX^{\mu_1}\wedge\cdots\wedge dX^{\mu_p}.
      \]
      Varying the embedding and using \(d^2=0\) gives
      \[
        \delta Z^{\mu_1\cdots\mu_p}
          =T_p\int_{\Sigma_p}d\!\left[
            \sum_{r=1}^p(-1)^{r-1}\delta X^{\mu_r}
            dX^{\mu_1}\wedge\cdots\widehat{dX^{\mu_r}}
            \cdots\wedge dX^{\mu_p}\right].
      \]
      Stokes' theorem turns this into an integral over
      \(\partial\Sigma_p\).  It therefore vanishes for a closed brane or
      for variations fixed at infinity.  The charge depends only on the
      homology or asymptotic winding class, not on local deformations.
      It transforms as a spacetime \(p\)-form and may appear in the
      supercharge anticommutator when the appropriate gamma-matrix
      bilinear is symmetric.
    \item \emph{Completeness of the \(p\)-Form Basis.}
    For a spinor space of dimension \(D_s\), the matrices
      \[
        {\bf1},\quad\gamma^\mu,\quad
        \gamma^{[\mu\nu]},\quad\ldots
      \]
      are mutually orthogonal under the trace pairing when their ranks
      differ.  Their formal total number is
      \[
        \sum_{p=0}^d\binom dp=2^d,
      \]
      which matches the matrix algebra of a Dirac spinor in even \(d\).
      In odd \(d\), rank \(p\) and rank \(d-p\) are Hodge-dual matrices, so
      only half are independent: \(2^{d-1}\), exactly the dimension of the
      matrix algebra of an irreducible \(2^{(d-1)/2}\)-component spinor.
      Hence the independent antisymmetrized products span every spinor
      matrix.

      Multiplication by the volume gamma matrix relates rank \(p\) to rank
      \(d-p\); on an irreducible chiral space one retains only one member
      of each dual pair.  Finally,
      \(\{Q_\alpha,Q_\beta\}\) is symmetric in \(\alpha,\beta\).
      Expanding it in the complete gamma basis therefore keeps exactly
      those \(\gamma^{[p]}C\) that are symmetric, with internal-index
      antisymmetry compensating when extended supersymmetry is present.
      No additional bosonic spinor structure is possible.
    \item \emph{The \(528\) Components of the M-Algebra.}
    A symmetric real \(32\times32\) matrix has
      \[
        \frac{32\cdot33}{2}=528
      \]
      independent entries.  In eleven dimensions the symmetric
      gamma-matrix bilinears have ranks \(p=1,2,5\).  Their component
      counts are
      \[
        \binom{11}{1}=11,\qquad
        \binom{11}{2}=55,\qquad
        \binom{11}{5}=462.
      \]
      Since
      \[
        11+55+462=528,
      \]
      they exhaust the symmetric anticommutator:
      \[
        \{Q,Q\}\sim
          \Gamma^M P_M+\Gamma^{MN}Z_{MN}
           +\Gamma^{M_1\cdots M_5}Z_{M_1\cdots M_5}.
      \]
      The two- and five-form terms are the M2- and M5-brane charges.
  \end{enumerate}

  \SupplementarySolution{9}{Self-Dual Tensor Fields and Their Branes}
  \begin{enumerate}
    \item On a \(p\)-form in a spacetime with one timelike direction,
      \[
        *^2=(-1)^{p(D-p)+1}.
      \]
      For a middle form, \(p=n=D/2\), this becomes
      \[
        *^2=(-1)^{n^2+1}.
      \]
      A nonzero real eigenform satisfying \(*F_n=\pm F_n\) requires
      \(*^2=1\), hence \(n\) must be odd.  Writing \(n=2k+1\) gives
      \[
        D=2n=4k+2.
      \]

    \item The field strength \(F_n=dA_{n-1}\) belongs before imposing
      self-duality to a massless \((n-1)\)-form gauge potential.  Its
      little-group polarization count is
      \[
        N_{\rm unconstrained}=\binom{D-2}{n-1}.
      \]
      At middle degree the Hodge star splits the physical polarizations
      into two equal eigenspaces.  A self-dual or anti-self-dual field has
      \[
        N_{\rm SD}=\frac12\binom{D-2}{n-1}
          =\frac12\binom{D-2}{D/2-1}.
      \]
      For example, a chiral two-form in \(D=6\) has
      \(\tfrac12\binom42=3\) polarizations.

    \item An electric \(p\)-brane couples through
      \(\int_{\Sigma_{p+1}}A_{p+1}\).  Here the potential has rank
      \(n-1\), so
      \[
        p+1=n-1,
        \qquad p=n-2=\frac D2-2=2k-1.
      \]
      The magnetic brane associated with an \((n-1)\)-form potential has
      spatial dimension \(D-(n-1)-3=D-n-2=n-2\), the same answer.  This
      coincidence is the brane-level reflection of middle-degree
      self-duality.

    \item With no timelike direction, the Hodge formula is instead
      \[
        *^2=(-1)^{p(D-p)}.
      \]
      On a middle form it is \((-1)^{n^2}\).  Real eigenvalues \(\pm1\)
      therefore occur when \(n\) is even, namely in \(D=4k\).  In
      Euclidean \(D=4k+2\), \(*^2=-1\); after complexification the
      eigenvalues are \(\pm i\), so the Lorentzian real self-duality
      condition becomes an imaginary self-duality condition rather than a
      real one.
  \end{enumerate}

  \SupplementarySolution{10}{Dp-Brane Spectra and Type-II Parity}
  \begin{enumerate}
    \item \emph{Dp-Brane Worldvolume Degrees of Freedom.}
    A massless gauge field in \(p+1\) dimensions has
      \[
        (p+1)-2=p-1
      \]
      physical polarizations when \(p\geq1\).  A Dp-brane also breaks
      translations in the \(9-p\) transverse spatial directions, producing
      one scalar for each broken translation.  Hence
      \[
        n_B=(p-1)+(9-p)=8.
      \]
      This is also the dimensional reduction of the eight physical
      polarizations of a ten-dimensional super-Yang--Mills vector.

      The formula \(D-2\) for a massless vector assumes \(D\geq2\), so it
      cannot be continued literally to \(p=0\).  On a worldline \(A_0\) is
      a Lagrange multiplier and the D0 transverse coordinates are quantum
      mechanical variables rather than propagating massless waves.  Their
      counting must instead be done in the constrained Hamiltonian system;
      it does not contradict the light-cone open-string count used for
      \(p\geq1\).
    \item \emph{Even and Odd D-Branes.}
    A Dp-brane preserves supersymmetry parameters satisfying a condition
      of the schematic form
      \[
        \epsilon_2=\Gamma^{0\cdots p}{\cal K}_p\epsilon_1,
      \]
      where \({\cal K}_p\) is the appropriate chirality/Pauli-matrix
      factor.  A product of \(p+1\) gamma matrices reverses
      ten-dimensional chirality when \(p+1\) is odd and preserves it when
      \(p+1\) is even.

      In IIA, \(\epsilon_1\) and \(\epsilon_2\) have opposite chirality, so
      the projector requires \(p+1\) odd, or \(p\) even.  In IIB they have
      the same chirality, so \(p+1\) must be even, or \(p\) odd.  Thus
      \[
        \text{IIA}:p=0,2,4,6,8,\qquad
        \text{IIB}:p=-1,1,3,5,7,9,
      \]
      subject to the usual interpretation of the instantonic and
      space-filling cases.
  \end{enumerate}

  \SupplementarySolution{11}{Open-String Vector Multiplets and No Force}
  \begin{enumerate}
    \item \emph{BPS Tension, Charge, and No Force.}
    In the brane rest frame, positivity of the supercharge
      anticommutator gives an energy-density bound
      \[
        T\geq |q|.
      \]
      A supersymmetric brane has a nonzero spinor annihilated by the
      corresponding projector, so an eigenvalue of the anticommutator
      vanishes.  The inequality is saturated:
      \(T=|q|\), after choosing the common normalization of the
      gravitational and Ramond--Ramond kinetic terms.

      The long-range exchange between two parallel identical branes has
      two parts.  Graviton and dilaton exchange is attractive and
      proportional to \(T^2\); exchange of the Ramond--Ramond
      \((p+1)\)-form between like charges is repulsive and proportional to
      \(q^2\).  BPS equality makes the magnitudes equal, so the static force
      cancels.  Reversing one charge produces a brane--antibrane pair; the
      form-field force then becomes attractive and the cancellation is
      lost.
    \item \emph{The Open-String Vector Multiplet.}
    The Ramond ground state is a ten-dimensional spinor.  The GSO
      projection keeps one Majorana-Weyl chirality, with sixteen real
      components.  The massless Dirac equation and its associated
      light-cone constraint halve this to eight physical fermionic
      polarizations.

      The Neveu--Schwarz sector supplies the eight bosonic modes counted in
      Solution~25.  Hence the massless spectrum has
      \[
        8_B+8_F
      \]
      states.  It is the dimensional reduction to \(p+1\) dimensions of
      the ten-dimensional \(N=1\) super-Yang--Mills multiplet: a gauge
      field, \(9-p\) adjoint scalars, and their gaugino partners.  For one
      brane the gauge group is \(U(1)\); coincident branes promote it to a
      nonabelian \(U(N)\) multiplet.
  \end{enumerate}

  \SupplementarySolution{12}{T-Duality and Threshold Bounds}
  \begin{enumerate}
    \item \emph{T-Duality Changes Brane Dimension.}
    T-duality interchanges
      \(\partial_\tau X\) and \(\partial_\sigma\widetilde X\).  Therefore a
      Neumann condition
      \(\partial_\sigma X=0\) becomes a Dirichlet condition
      \(\partial_\tau\widetilde X=0\), while a Dirichlet condition becomes
      Neumann.

      Dualizing a direction parallel to a Dp-brane converts one Neumann
      coordinate to Dirichlet.  The brane loses that worldvolume direction
      and becomes a D\((p-1)\)-brane.  Dualizing a transverse direction
      converts one Dirichlet coordinate to Neumann, so the brane gains a
      worldvolume direction and becomes a D\((p+1)\)-brane.  Because one
      spacetime chirality is reversed at the same time, this map exchanges
      the even-\(p\) IIA spectrum with the odd-\(p\) IIB spectrum.
    \item \emph{Brane Threshold Bound.}
    In the rest frame, write the central-charge contribution as a
      Hermitian matrix \({\cal Z}\) on the supercharge space.  A unitary
      transformation diagonalizes it, and positivity of
      \[
        \{Q,Q^\dagger\}\sim M{\bf1}-{\cal Z}
      \]
      requires
      \[
        M\geq\lambda_{\max}({\cal Z})
      \]
      together with the corresponding inequality for the opposite signs.
      Equivalently, \(M\) is no smaller than the largest absolute charge
      eigenvalue.  Saturation makes the associated eigenspace of
      supercharges annihilate the state and determines the preserved
      fraction.

      If compatible constituent projectors make the largest eigenvalue the
      sum of the separate charge magnitudes, then
      \(M_{\rm BPS}=\sum_i M_i\).  The constituents may separate at no
      energy cost: this is a threshold or marginal bound state.  For
      charges entering a combined eigenvalue such as
      \(\sqrt{q_1^2+q_2^2}\), the saturated mass is smaller than
      \(|q_1|+|q_2|\).  The positive difference is binding energy, and the
      configuration is a nonthreshold BPS bound state.
  \end{enumerate}

\end{enumerate}
